\section{IBM Quantum Hardware Validation}\label{sec:hw}
\subsection{Scope and provenance}
Selected QSFD circuits were executed on IBM Quantum hardware to measure how far real-device behaviour departs from ideal and backend-derived noisy simulation. The campaign used \HwBackend{} with \HwJobs{} independent jobs of 52 circuits at \HwShots{} shots each (\HwCircuits{} circuits, \HwTotalShots{} shots), after a pilot job (\HwPilotJob{}; 17 circuits, 1024 shots) that is retained as a baseline. An earlier job on \texttt{ibm\_fez} (\HwFezJob{}) executed the Layer~1 QNN circuits and is reported separately. These runs validate circuit-level components; they do not execute the full fraud classifier and are not physical quantum key distribution.

\subsection{BB84 primitives and the intercept--resend model}
Table~\ref{tab:hw-backend} compares theory, ideal Aer, noisy Aer and hardware. With no Eve, hardware QBER was \HwNoEveBest{} (95\% CI \HwNoEveCIBest{}) on the best-calibrated qubit (physical qubit \HwQubitBest{}, readout error \HwReadoutBest{}), \HwNoEveTypical{} on the typical qubit and \HwNoEveWorse{} on the worse qubit (readout error \HwReadoutWorse{}); the ordering matches the readout-error ordering, although with three qubits only a descriptive comparison is made. The intercept--resend circuit model raised QBER to \HwEveBest{}, \HwEveTypical{} and \HwEveWorse{}, i.e.\ by \HwDeltaBest{}, \HwDeltaTypical{} and \HwDeltaWorse{} relative to no Eve (paired bootstrap CI for the best qubit [\HwDeltaBootLoBest{}, \HwDeltaBootHiBest{}]; exact sign-flip permutation $p=\HwPermBest{}$, the minimum attainable with five pairs). Attack-induced error stayed separable from hardware error in every job and qubit cell: the largest no-Eve QBER (\HwMaxNoEve{}) was below the smallest Eve QBER (\HwMinEve{}). The ideal expectation of 25\% is a reference rather than a pass/fail target; hardware values above it are consistent with hardware-induced error added to attack-induced error.

\begin{table}[t]
\centering
\small
\caption{Theory vs.\ ideal Aer vs.\ backend-derived noisy Aer vs.\ IBM hardware (ibm\_marrakesh); mean and 95\% $t$-interval over independent jobs.}
\label{tab:hw-backend}
\resizebox{\linewidth}{!}{%
\begin{tabular}{llrrrrrrr}
\toprule
Experiment & Qubit & Theory & Ideal Aer & Noisy Aer & IBM mean & CI low & CI high & Jobs \\
\midrule
BB84 matched-basis QBER (no Eve) & best (q98) & 0.0000 & 0.0000 & 0.0025 & 0.0048 & 0.0036 & 0.0060 & 5 \\
Intercept-resend QBER & best (q98) & 0.2500 & 0.2500 & 0.2512 & 0.2632 & 0.2579 & 0.2685 & 5 \\
BB84 matched-basis QBER (no Eve) & typical (q155) & 0.0000 & 0.0000 & 0.0110 & 0.0094 & 0.0082 & 0.0106 & 5 \\
Intercept-resend QBER & typical (q155) & 0.2500 & 0.2500 & 0.2557 & 0.2591 & 0.2535 & 0.2648 & 5 \\
BB84 matched-basis QBER (no Eve) & worse (q126) & 0.0000 & 0.0000 & 0.0738 & 0.0446 & 0.0277 & 0.0616 & 5 \\
Intercept-resend QBER & worse (q126) & 0.2500 & 0.2500 & 0.2873 & 0.2868 & 0.2796 & 0.2939 & 5 \\
Grover 2q success probability & pair [135, 139] & 1.0000 & 1.0000 & 0.9678 & 0.9460 & 0.9386 & 0.9534 & 5 \\
\bottomrule
\end{tabular}}
\end{table}

\begin{table}[t]
\centering
\small
\caption{Physical qubits used and their calibration snapshot at campaign start, with observed hardware error.}
\label{tab:hw-calibration}
\resizebox{\linewidth}{!}{%
\begin{tabular}{lrrrrrlr}
\toprule
Role & Qubit & T1 (\$\textbackslash mu\$s) & T2 (\$\textbackslash mu\$s) & Readout err. & 1Q err. & Circuit & No-Eve QBER \\
\midrule
best & 98 & 265.6 & 68.85 & 0.002075 & 0.0008155 & BB84 single-qubit & 0.004785 \\
typical & 155 & 393.9 & 107 & 0.01221 & 0.0003537 & BB84 single-qubit & 0.009399 \\
worse & 126 & 104.3 & 8.646 & 0.07593 & 0.0004358 & BB84 single-qubit & 0.04463 \\
grover\_pair & 135 & 294.5 & 91.06 & 0.02295 & 0.0001081 & Grover 2q (cz err 0.0011) & -- \\
grover\_pair & 139 & 136.3 & 158.9 & 0.005371 & 0.00019 & Grover 2q (cz err 0.0011) & -- \\
\bottomrule
\end{tabular}}
\end{table}


\begin{figure}[t]
\centering
\includegraphics[width=\figw]{hw_qber_theory_ideal_noisy_hardware}
\caption{BB84 QBER: theory, ideal Aer, backend-derived noisy Aer and IBM hardware (95\% CI over independent jobs).}
\label{fig:hwqber}
\end{figure}
\begin{figure}[t]
\centering
\includegraphics[width=\figw]{hw_qber_per_job}
\caption{Run-to-run variability of hardware no-Eve QBER: one point per independent job.}
\label{fig:hwjobs}
\end{figure}
\begin{figure}[t]
\centering
\includegraphics[width=\figw]{hw_eve_vs_noeve}
\caption{Hardware no-Eve versus intercept--resend-model QBER by job and qubit tier.}
\label{fig:hweve}
\end{figure}

\subsection{State asymmetry and simulator prediction}
The pilot suggested larger error for $|1\rangle$-prepared states. In the replicated data, $P(\text{misread}\mid |1\rangle)$ exceeded $P(\text{misread}\mid |0\rangle)$ in \HwAsymCells{} job-by-qubit cells: on the best qubit the values were \HwAsymZeroBest{} for $|0\rangle$ and \HwAsymOneBest{} for $|1\rangle$, but on the worse qubit the pattern reversed (\HwAsymZeroWorse{} versus \HwAsymOneWorse{}). The asymmetry is therefore qubit-dependent; it is consistent with relaxation and asymmetric readout effects, but no unique cause is established. Backend-derived noisy Aer predicted hardware QBER with MAE \HwNoisyMAE{} and RMSE \HwNoisyRMSE{} over 36 circuit-by-qubit cells (ideal Aer: MAE \HwIdealMAE{}); it over-predicted error on the worse qubit, and it used backend properties at analysis time rather than at job time.

\begin{figure}[t]
\centering
\includegraphics[width=\figw]{hw_state_asymmetry}
\caption{State-dependent misreads by qubit tier, with backend readout error (dashed).}
\label{fig:hwasym}
\end{figure}
\begin{figure}[t]
\centering
\includegraphics[width=\figw]{hw_noisy_aer_vs_hardware}
\caption{Noisy-Aer prediction versus hardware QBER per circuit and qubit tier.}
\label{fig:hwnoisy}
\end{figure}

\subsection{Grover and transpilation}
A two-qubit Grover search over four items with marked state $|11\rangle$ (ideal success 1) reached a pooled hardware success probability of \HwGroverPooled{}. By optimization level 0 to 3 the hardware success was \HwGroverOptZero{}, \HwGroverOptOne{}, \HwGroverOptTwo{} and \HwGroverOptThree{} (noisy Aer about \HwGroverNoisyOne{}), with transpiled depths \HwGroverDepthZero{}, \HwGroverDepthOne{}, \HwGroverDepthTwo{}, \HwGroverDepthThree{} and \HwGroverTwoQOne{} two-qubit gates at every level; differences between levels are small relative to between-job variation, so no optimization-level trend is claimed. This is an educational demonstration with no implication for AES or RSA key sizes; Shor's algorithm was not implemented.

\begin{table}[t]
\centering
\small
\caption{Two-qubit Grover (marked state $|11\rangle$): success probability by transpiler optimization level (5 independent jobs).}
\label{tab:hw-grover}
\resizebox{\linewidth}{!}{%
\begin{tabular}{rrrrrrrr}
\toprule
Opt. level & Depth & 2Q gates & HW success & CI low & CI high & Noisy Aer & Ideal Aer \\
\midrule
0 & 14 & 2 & 0.9465 & 0.9357 & 0.9573 & 0.9673 & 1.0000 \\
1 & 12 & 2 & 0.9521 & 0.9454 & 0.9589 & 0.9680 & 1.0000 \\
2 & 10 & 2 & 0.9441 & 0.9303 & 0.9580 & 0.9680 & 1.0000 \\
3 & 10 & 2 & 0.9412 & 0.9318 & 0.9506 & 0.9680 & 1.0000 \\
\bottomrule
\end{tabular}}
\end{table}

\begin{figure}[t]
\centering
\includegraphics[width=\figw]{hw_grover_success}
\caption{Two-qubit Grover success probability: ideal Aer, noisy Aer and hardware.}
\label{fig:hwgrover}
\end{figure}

\subsection{QML circuits on \texttt{ibm\_fez}}
The CLT-QNN and linear-QNN circuits (4 qubits) with simulator-trained, frozen weights were bound to six held-out transactions and executed once. Mean absolute deviation between hardware and simulator expectation values was \FezCltMAD{} (Chow--Liu, transpiled depth \FezCltDepth{}) and \FezLinearMAD{} (linear, depth \FezLinearDepth{}), with hardware outputs strongly compressed toward zero for the deeper linear circuit. This supports only the limited claim that the circuits are physically executable and that hardware outputs differ markedly from simulation (more so for the deeper circuit; two circuits only); it is not evidence of a hardware predictive advantage, and the training experiment itself is simulator-based.
